\section{Start with the paper's question}
\label{sec:start}

A paper becomes difficult to write when its authors have not fixed its central
question and claim. Polishing cannot settle either one. Before drafting, state
the scientific problem and the paper's answer in ordinary language.

\subsection{Write the paper's contract}

Write four sentences before writing any section:

\begin{enumerate}[leftmargin=*,label=\arabic*.]
    \item \textbf{Problem.} What practical or scientific problem does the paper
    address?
    \item \textbf{Gap.} Why does the strongest current method or explanation
    leave that problem unresolved?
    \item \textbf{Claim.} What does this paper establish?
    \item \textbf{Boundary.} Where does the claim stop, and what important
    question remains open?
\end{enumerate}

These sentences form a contract with the reader. The introduction may motivate
the problem and the body may develop the details, but neither should quietly
change the contract. The boundary sentence matters as much as the central
claim. It prevents a result in one regime from becoming a claim about every
regime. It also keeps a measured effect separate from the explanation proposed
for it.

A useful one-sentence test is:

\begin{quote}
We study \emph{[problem]} because \emph{[current limitation]}; we show
\emph{[central result]} under \emph{[scope]}, while \emph{[limitation]} remains
open.
\end{quote}

The finished sentence need not appear in the paper. Its purpose is to expose
disagreement among the authors. If two plausible versions imply different
experiments or different section orders, settle that disagreement before
line-editing the introduction.

Match the intended audience to the reach of the evidence. State the new
understanding the paper offers that audience, then look to adjacent fields for
a useful framing. That comparison can strengthen the argument without
inflating the claim.

\subsection{Map claims to evidence}

List every claim required for the paper's conclusion to hold.
Then assign each claim a visible piece of evidence. A theorem supports a formal
statement under its assumptions. A controlled comparison supports a
comparative claim. An ablation isolates a design choice. A mechanism study
tests an explanation. A qualitative result can make a failure mode concrete,
but it rarely establishes frequency on its own.

\begin{center}
\begin{tabular}{p{0.20\linewidth}p{0.31\linewidth}p{0.31\linewidth}}
\toprule
\textbf{Claim} & \textbf{Evidence needed} & \textbf{Boundary to state} \\
\midrule
Headline result &
Theorem or controlled benchmark &
Assumptions, tasks, metrics, and uncertainty \\
\addlinespace
Mechanism &
Diagnostic intervention or measured intermediate behavior &
Alternative explanations that remain possible \\
\addlinespace
Practical value &
Cost, deployment result, or stress test &
Resources and settings not yet tested \\
\bottomrule
\end{tabular}
\end{center}

This map should reveal weak links. A contribution with no evidence is a
proposal, not a result. Evidence that does not support the central argument may
still belong in the appendix, but it should not dictate the main paper. When a
claim depends on a missing control, record the gap instead of drafting around
it.

\subsection{Choose the level of the contribution}

Papers make different kinds of contributions. A method paper may show that a
new procedure resolves a concrete tradeoff. A mechanism paper may overturn a
popular explanation through controlled tests. A theory paper may identify a
failure and prove when it disappears. Some papers combine these forms, but one
usually carries the main argument.

Defining a constrained regime can also be a contribution. A careful problem
definition may expose assumptions and design choices that abundant data or a
looser formulation hides. In that case, say that the paper defines the regime,
tests a recipe, or identifies necessary conditions. Do not say that it solves
the broader problem.

Name that primary contribution early. It determines the evidence that deserves
space and the content of Figure~1. It should also govern the title and
abstract. A paper centered on explanation should not market itself as an
algorithm paper because it happens to introduce a diagnostic model. A paper
centered on a theorem should not bury the regime in which the theorem applies.
